\documentclass[11pt]{article}
\usepackage[T1]{fontenc}
\usepackage[utf8]{inputenc}
\usepackage{amsmath,amsthm,mathtools}
\usepackage{newpxtext,newpxmath}
\usepackage[scaled=0.93]{sourcesanspro}
\usepackage[a4paper,margin=25mm,headheight=15pt,footskip=13mm]{geometry}
\usepackage{microtype}
\usepackage{booktabs,array,enumitem,longtable}
\usepackage{xcolor}
\usepackage[most]{tcolorbox}
\usepackage{fancyhdr,lastpage,titlesec}
\usepackage{hyperref}
\usepackage{listings}
\definecolor{Ink}{HTML}{16354A}
\definecolor{Teal}{HTML}{176B77}
\definecolor{Pale}{HTML}{EFF6F7}
\definecolor{Muted}{HTML}{566773}
\hypersetup{colorlinks=true,linkcolor=Teal,citecolor=Teal,urlcolor=Teal,
  pdftitle={A counterexample, exact spacing criteria, and a sharp smoothing threshold for products of q-integers},
  pdfauthor={Research note prepared with ChatGPT},pdfsubject={Unimodality; q-integers; exact counterexample; smoothing threshold}}
\pagestyle{fancy}
\fancyhf{}
\fancyhead[L]{\sffamily\small\color{Muted}Q-INTEGER UNIMODALITY}
\fancyhead[R]{\sffamily\small\color{Muted}Counterexample, exact criteria, smoothing threshold}
\fancyfoot[C]{\sffamily\small\color{Muted}\thepage\ / \pageref*{LastPage}}
\renewcommand{\headrulewidth}{0.3pt}
\titleformat{\section}{\Large\sffamily\bfseries\color{Ink}}{\thesection}{0.7em}{}
\titleformat{\subsection}{\large\sffamily\bfseries\color{Ink}}{\thesubsection}{0.7em}{}
\titlespacing*{\section}{0pt}{2.2ex plus .7ex minus .2ex}{1.1ex}
\titlespacing*{\subsection}{0pt}{1.7ex plus .5ex minus .2ex}{.8ex}
\newtheoremstyle{research}{8pt}{8pt}{\itshape}{}{
  \sffamily\bfseries\color{Ink}}{.}{.5em}{}
\theoremstyle{research}
\newtheorem{theorem}{Theorem}[section]
\newtheorem{proposition}[theorem]{Proposition}
\newtheorem{lemma}[theorem]{Lemma}
\newtheorem{corollary}[theorem]{Corollary}
\theoremstyle{definition}
\newtheorem{definition}[theorem]{Definition}
\newtheorem{remark}[theorem]{Remark}
\newcommand{\qi}[1]{[#1]_q}
\newcommand{\qir}[2]{[#1]_{q^{#2}}}
\newcommand{\coef}[2]{[q^{#1}]\,#2}
\newcommand{\SU}{\textup{SU}}
\newcommand{\Z}{\mathbb Z}
\newcommand{\R}{\mathbb R}
\newcommand{\floor}[1]{\left\lfloor #1\right\rfloor}
\newcommand{\TS}{\mathcal B}
\DeclareMathOperator{\supp}{supp}
\numberwithin{equation}{section}
\setlength{\parskip}{3.5pt}
\setlength{\parindent}{1em}
\setlist{itemsep=3pt,topsep=5pt}
\tcbset{resultbox/.style={enhanced,breakable,colback=Pale,colframe=Teal,
  boxrule=.6pt,arc=1.5mm,left=3mm,right=3mm,top=2mm,bottom=2mm,
  fonttitle=\sffamily\bfseries,coltitle=white}}
\lstset{basicstyle=\ttfamily\small,columns=fullflexible,keepspaces=true,
  breaklines=true,frame=single,rulecolor=\color{Muted},showstringspaces=false,
  xleftmargin=3pt,xrightmargin=3pt}
\emergencystretch=1.5em

\begin{document}
\thispagestyle{plain}
\begin{center}
  {\sffamily\small\bfseries\color{Teal}RESEARCH NOTE \quad / \quad 18--20 SEPTEMBER 2026}\par
  \vspace{7pt}
  {\sffamily\bfseries\color{Ink}\LARGE A counterexample, exact spacing criteria,\par
  and a sharp smoothing threshold\par
  for products of \boldmath$q$\unboldmath-integers}\par
  \vspace{8pt}
  {\small Prepared with ChatGPT; self-contained proofs and exact computational checks}
\end{center}

\begin{abstract}
We attack Conjecture 5.4 of Connelly, Ito, Martinez, Shevchenko, and Yang,
\emph{Unimodality of $q$-Fibonomial coefficients for small cases},
arXiv:2605.12822v1. Its necessity assertion is false:
$(1+q)^6(1+q^3)$ is unimodal but violates the proposed condition.
We then prove the conjectured sufficient condition for every spacing $r\ge2$.
For spacing three we obtain a complete replacement: if no $a_i$ is divisible
by three, then
\[
 \left(\prod_i\qi{a_i}\right)\qir{b}{3}\text{ is unimodal}
 \quad\Longleftrightarrow\quad
 b\le 1+\sum_i\left\lfloor\frac{a_i}{3}\right\rfloor
       +2\left\lfloor\frac{\#\{i:a_i\equiv2\pmod3\}}{6}\right\rfloor.
\]
A failed inequality produces a specified first-half coefficient drop of
exactly one. We also establish a general necessary degree bound and settle
the spacing-two criterion. A dedicated study of the two-shift family
$(1+q)^n(1+q^r)$ proves the sharp threshold ``unimodal if and only if
$n\ge r^2-3$'' in two independent ways, determines the equality case
(exactly four consecutive maximal coefficients at the threshold), shows that
the spacing-three counterexamples are log-concave for every $n\ge6$, and
records a strict-maximum witness $(1+q)^9(1+q^3)$. A complete finite search
over \emph{all} spacings proves that degree nine is the minimum possible
degree of a counterexample and that its parameter set is unique.
The arguments are finite and algebraic; computations are supplementary.
The main $q$-Fibonomial unimodality conjecture is not resolved here.
\end{abstract}

\begin{tcolorbox}[resultbox,title={The counterexample at a glance, and a strict companion}]
\[
 (1+q)^6(1+q^3)
 =1+6q+15q^2+21q^3+21q^4+21q^5+21q^6+15q^7+6q^8+q^9.
\]
Here $r=3$, $b=2$, and $a_1=\cdots=a_6=2$.
No $a_i$ is divisible by $3$, and
$1+\sum_i\lfloor a_i/3\rfloor=1<2$.
Nevertheless, the displayed coefficient sequence is weakly unimodal.
It is in fact log-concave, and three further factors $1+q$ give
\begin{align*}
 (1+q)^9(1+q^3)={}&1+9q+36q^2+85q^3+135q^4+162q^5+168q^6\\
   &+162q^7+135q^8+85q^9+36q^{10}+9q^{11}+q^{12},
\end{align*}
which has a \emph{unique} maximum and strictly monotone flanks while still
violating the proposed condition. The refutation is therefore not an artefact
of the central plateau. Both witnesses are carried throughout
(item~D5 of Table~\ref{tab:dual}).
\end{tcolorbox}

\noindent\textbf{Status.} The statements below are accompanied by proofs and
reproducible integer-arithmetic tests. They have not been independently
refereed or checked by a proof assistant. The counterexample refutes the
specified source statement; historical priority of the counterexample and
the stronger theorems has not been established by an exhaustive literature review.

\section{The selected conjecture and the result}
\label{sec:target}
For a positive integer $a$, put
\[
 \qi{a}=1+q+\cdots+q^{a-1}.
\]
Throughout, $q$ is a formal indeterminate, $r\ge2$, $k\ge1$, and
$a_1,\ldots,a_k,b$ are positive integers. We study
\begin{equation}\label{eq:product}
 P_{\boldsymbol a,b,r}(q)
   =\left(\prod_{i=1}^{k}\qi{a_i}\right)\qir{b}{r}.
\end{equation}
Its coefficients are nonnegative integers, its constant and leading
coefficients are one, and its degree is
\begin{equation}\label{eq:degree}
 D=S+r(b-1),\qquad S=\sum_{i=1}^k(a_i-1).
\end{equation}
Every factor is symmetric, so $P$ is symmetric.

The source of the attack is \cite[Section 5.2, Conjecture 5.4, p.~14]{CIMSY}.
For brevity write its proposed condition as
\begin{equation}\label{eq:condition}
 \mathcal C_r(\boldsymbol a,b):\quad
 (\exists i\colon r\mid a_i)
 \quad\text{or}\quad
 b\le 1+\sum_i\left\lfloor\frac{a_i}{r}\right\rfloor.
\end{equation}
The conjecture asserts that \eqref{eq:condition} is sufficient, and also
necessary when $k\le3$ \emph{or} $r\le3$.
The distinction between ``or'' and ``and'' matters: the case
$k=6,r=3$ lies inside the asserted necessity range.
The source reports tests with $k\le5$, $r\le6$, and
$\max\{a_i,b\}\le15$ \cite[p.~14]{CIMSY}.
This is an auxiliary conjecture inside a paper about $q$-Fibonomial
coefficients; it is not the original Fibonomial unimodality conjecture
of \cite{BCK}.

\begin{proposition}[Exact disproof]\label{prop:counterexample}
The necessity assertion in \cite[Conjecture 5.4]{CIMSY} is false.
\end{proposition}
\begin{proof}
Take $r=3$, $b=2$, and $a_i=2$ for $1\le i\le6$.
The binomial theorem gives the coefficient list
\[
 (1,6,15,21,21,21,21,15,6,1).
\]
It increases weakly to its central plateau and then decreases weakly.
On the other hand, $3\nmid2$ and
$b=2>1+6\lfloor2/3\rfloor=1$, so \eqref{eq:condition} fails.
Since $r=3$, the claimed necessity applies to this example.
\end{proof}

The complete local certificate is small enough to print in full:
\begin{center}
\small
\begin{tabular}{c@{\quad}rrrrrrrrrr}
\toprule
$j$ &0&1&2&3&4&5&6&7&8&9\\
\midrule
$\coef jP$&1&6&15&21&21&21&21&15&6&1\\
Difference from previous&---&5&9&6&0&0&0&$-6$&$-9$&$-5$\\
\bottomrule
\end{tabular}
\end{center}

This is a counterexample to an auxiliary product criterion, not to the
main $q$-Fibonomial conjecture in \cite{CIMSY}.
Our stronger conclusions, proved below, are summarized in
Table~\ref{tab:results}. All appearances of ``unimodal'' allow equal
consecutive coefficients; a sequence is unimodal if it is weakly increasing
up to some index and weakly decreasing thereafter, intervening zeros included.

\begin{table}[htbp]
\centering
\small
\begin{tabular}{@{}p{.25\linewidth}p{.64\linewidth}@{}}
\toprule
Setting & Conclusion proved in this note \\
\midrule
Every $r\ge2$ & The condition $\mathcal C_r$ is sufficient
(Theorem~\ref{thm:sufficient}). With no divisible
factor, unimodality also requires $b\le1+\lfloor S/r\rfloor$
(Theorem~\ref{thm:degreebound}). \\
$r=2$ & The original condition $\mathcal C_2$ is necessary and sufficient
(Theorem~\ref{thm:r2}). \\
$r=3$ & With no divisible factor, the exact upper bound on $b$ is
$1+\sum_i\lfloor a_i/3\rfloor+2\lfloor t/6\rfloor$,
where $t=\#\{i:a_i\equiv2\pmod3\}$ (Theorem~\ref{thm:r3}). \\
Two shifted binomial rows & $(1+q)^n(1+q^r)$ is unimodal exactly when
$n\ge r^2-3$, for $r\ge2$, $n\ge0$ (Theorem~\ref{thm:twoshift}, with two
independent proofs). At the threshold exactly four consecutive coefficients
are maximal (Corollary~\ref{cor:plateau}). \\
Gap three, all $n\ge6$ & $(1+q)^n(1+q^3)$ is a \emph{log-concave}
counterexample (Corollary~\ref{cor:infinite}); and for every $B\ge2$,
$(1+q)^t\qir B3$ with $t\ge6\lfloor B/2\rfloor$ is one as well
(Corollary~\ref{cor:binomialexact}). \\
All spacings & Minimum counterexample degree is nine, with a unique
normalized parameter set (Proposition~\ref{prop:minimum}). \\
\bottomrule
\end{tabular}
\caption{Mathematical results, rather than extrapolations from the finite tests.}
\label{tab:results}
\end{table}

\paragraph{How this note is organized.}
Sections~\ref{sec:tools}--\ref{sec:r3} treat the general product
\eqref{eq:product}: closure tools, the centered peeling identity and
sufficiency, a residue-mass obstruction and a necessary degree bound, and
complete criteria for spacings two and three.
Section~\ref{sec:twoshift} is a deep dive on the two-shift slice
$b=2$, $a_i=2$, which is exactly the family containing the counterexample.
Section~\ref{sec:minimal} settles minimum degree across all spacings,
Section~\ref{sec:probability} gives a clearly labelled probabilistic
heuristic for the quadratic scale, Section~\ref{sec:verification} records
the computations, Section~\ref{sec:provenance} the provenance, and
Section~\ref{sec:scope} what remains open.

\section{Elementary tools for symmetric unimodality}
\label{sec:tools}
A finite nonnegative sequence $(c_0,\ldots,c_D)$ is \emph{unimodal} if,
for some index $m$, it is weakly increasing up to $m$ and weakly decreasing
after $m$. It is \emph{symmetric} if $c_j=c_{D-j}$.
For a symmetric sequence, unimodality is equivalent to
\begin{equation}\label{eq:firsthalf}
 c_j-c_{j-1}\ge0\qquad
 0\le j\le\lfloor D/2\rfloor,
 \qquad c_{-1}=0.
\end{equation}
Indeed, these inequalities control the left half and reflection controls the
right half; in odd degree the two central entries agree automatically.
We abbreviate ``symmetric and unimodal'' by $\SU$.
All coefficient sequences are extended by zero outside their support.

A shifted sequence $q^sH(q)$ has its center shifted by $s$.
In a sum of shifted polynomials, symmetry is understood about the
\emph{specified common center}, not automatically half the largest
exponent of each summand. This convention is essential in the decomposition
used below.

\begin{lemma}[Centered intervals and closure]\label{lem:closure}
Products of $\SU$ polynomials with nonnegative coefficients are $\SU$.
Nonnegative sums of $\SU$ coefficient sequences with the same center are
also $\SU$.
\end{lemma}
\begin{proof}
If $F(q)=\sum_{j=0}^D f_jq^j$ is $\SU$, let
$\alpha_j=f_j-f_{j-1}$ for $0\le j\le\lfloor D/2\rfloor$, with
$f_{-1}=0$. Then $\alpha_j\ge0$ and
\begin{equation}\label{eq:intervals}
 F(q)=\sum_{j=0}^{\lfloor D/2\rfloor}
          \alpha_j q^j\qi{D-2j+1}.
\end{equation}
To verify the identity, the coefficient at $q^m$ on the right is
$\sum_{j=0}^{\min(m,D-m)}\alpha_j=f_{\min(m,D-m)}=f_m$.
Each summand is the indicator sequence of an interval centered at $D/2$.

For positive $u,v$, the coefficients of $\qi{u}\qi{v}$ at exponents
$0\le\ell\le u+v-2$ are
\[
 \min\{\ell+1,u,v,u+v-1-\ell\}.
\]
They form a symmetric triangular or trapezoidal sequence.
Apply \eqref{eq:intervals} to two factors of degrees $D$ and $E$.
Every product of two interval summands is $\SU$ with center $(D+E)/2$.
On the left of this center, each first difference is nonnegative, so the
same is true of their sum. Symmetry supplies the right half. This also
proves the assertion about sums with a common center.
\end{proof}

The special case in which one factor is $1+q$ admits a two-line proof that
avoids the interval decomposition entirely. We record it separately because
it is the only closure fact the two-shift analysis of
Section~\ref{sec:twoshift} needs, and because it also yields the plateau
bookkeeping of Section~\ref{sec:twoshift-equality}.

\begin{lemma}[$(1+q)$-smoothing]\label{lem:smoothing}
If $A(q)$ has nonnegative, symmetric, unimodal coefficients, then so does
$(1+q)A(q)$.
\end{lemma}
\begin{proof}
Write $A(q)=\sum_{j=0}^{D}a_jq^j$, extending $a_j$ by zero outside its
support. The new coefficients satisfy
\begin{equation}\label{eq:smoothing-difference}
 b_j=a_j+a_{j-1},\qquad b_j-b_{j-1}=a_j-a_{j-2}.
\end{equation}
Symmetry and nonnegativity are immediate. For $j$ in the left half of the
new sequence, $a_j\ge a_{j-2}$ follows from the old left-half monotonicity,
except possibly at the new middle index when $D=2h+1$. At that index
$j=h+1$, and symmetry gives $a_{h+1}=a_h\ge a_{h-1}$, which is the required
inequality. Now apply \eqref{eq:firsthalf}.
\end{proof}

Lemmas~\ref{lem:closure} and \ref{lem:smoothing} are kept as a deliberate
pair; see item~D10 of Table~\ref{tab:dual}.

\begin{lemma}[A divisible factor]\label{lem:divisible}
If $r\mid a_i$ for some $i$, then $P_{\boldsymbol a,b,r}$ is unimodal for
every positive integer $b$.
\end{lemma}
\begin{proof}
Write $a_i=rA$. The identity
\[
 \qi{rA}=\qi r\qir A r
\]
shows that the relevant two factors are
\[
 \qi r\qir A r\qir b r.
\]
The polynomial $[A]_z[b]_z$ has an $\SU$ coefficient sequence in $z$.
Replacing $z$ by $q^r$ and multiplying by $\qi r$ repeats each coefficient
exactly $r$ times. Repetition preserves symmetry and unimodality.
The remaining factors are ordinary $q$-integers, so
Lemma~\ref{lem:closure} finishes the proof.
\end{proof}

\section{A centered identity proves the sufficient condition}
\label{sec:sufficiency}
The following identity removes one block of length $r$ from an ordinary
factor and one term from the spaced factor, without changing the center
of the residual sequence.

\begin{lemma}[Centered peeling]\label{lem:peeling}
For integers $a>r$ and $b\ge2$,
\begin{equation}\label{eq:peeling}
 \qi a\qir b r
 =\qi{a+r(b-1)}+q^r\qi{a-r}\qir{b-1}{r}.
\end{equation}
Both terms on the right, with zero padding, are symmetric about the
center of the polynomial on the left.
\end{lemma}
\begin{proof}
Multiply \eqref{eq:peeling} by $(1-q)(1-q^r)$. The identity becomes
\begin{align*}
 (1-q^a)(1-q^{rb})
 &= (1-q^{a+r(b-1)})(1-q^r)\\
 &\quad +q^r(1-q^{a-r})(1-q^{r(b-1)}),
\end{align*}
which follows on expansion. The left side of \eqref{eq:peeling} has
polynomial degree
\[
 d=a-1+r(b-1).
\]
The first term on the right has degree $d$. The unshifted residual
polynomial has degree $d-2r$, so its shift by $q^r$ has center
$r+(d-2r)/2=d/2$.
\end{proof}

\begin{theorem}[Sufficiency for every spacing]\label{thm:sufficient}
For every $r\ge2$ and every number $k$ of ordinary factors,
$\mathcal C_r(\boldsymbol a,b)$ implies unimodality of
$P_{\boldsymbol a,b,r}$.
\end{theorem}
\begin{proof}
A divisible factor is covered by Lemma~\ref{lem:divisible}. Assume there
is none, and write
\[
 Q=\sum_i\left\lfloor\frac{a_i}{r}\right\rfloor.
\]
We prove sufficiency of $b\le Q+1$ by induction on $b$.
For $b=1$, the spaced factor equals one and the product of ordinary
$q$-integers is $\SU$ by Lemma~\ref{lem:closure}.

Let $b\ge2$. Then $Q\ge1$, so at least one $a_i$ is greater than $r$:
equality is impossible because there is no divisible factor.
Apply \eqref{eq:peeling} to this factor and multiply by all the others.
The first summand is a product of ordinary $q$-integers, hence $\SU$.
In the residual product, $a_i$ is replaced by $a_i-r$ and $b$ by $b-1$.
No new divisible factor appears, the floor sum becomes $Q-1$, and
\[
 b-1\le Q=(Q-1)+1.
\]
The induction hypothesis makes the residual product $\SU$.
Lemma~\ref{lem:peeling} shows that the shifted residual and the first
summand have the same center, even after multiplication by the other
factors. Their sum is therefore $\SU$ by Lemma~\ref{lem:closure}.
\end{proof}

This settles the half of \cite[Conjecture 5.4]{CIMSY} that the
counterexample does not touch. It is proved here, not merely tested; the
degree-nine exhaustive search of Section~\ref{sec:minimal} independently
found no violation of sufficiency, which is a corroborating bounded
observation rather than evidence for the theorem.

\begin{corollary}[Reduction to small residues]\label{cor:core}
Suppose no $a_i$ is divisible by $r$, and write
\[
 a_i=rA_i+s_i,\quad 1\le s_i\le r-1,\qquad Q=\sum_i A_i.
\]
If $b>Q+1$ and
\begin{equation}\label{eq:core}
 \left(\prod_i\qi{s_i}\right)\qir{b-Q}{r}
\end{equation}
is $\SU$, then $P_{\boldsymbol a,b,r}$ is $\SU$.
\end{corollary}
\begin{proof}
Apply \eqref{eq:peeling} a total of $Q$ times, subtracting $r$ from
some $a_i$ at each step. The spaced parameter never drops below two,
since its final value is $b-Q\ge2$. The final residual is
\eqref{eq:core}. Reversing the steps, each step adds an $\SU$ product of
ordinary $q$-integers at the same center as the shifted residual.
\end{proof}

Corollary~\ref{cor:core} is a sufficient reduction, not an assertion that
residue reduction preserves unimodality in both directions for every $r$.
That converse will be obtained for $r=3$ by a separate obstruction argument.

\section{Periodic coefficient differences and a necessary bound}
\label{sec:residues}
Let
\[
 F(q)=\prod_i\qi{a_i}=\sum_{j=0}^S f_jq^j,
 \qquad R_s=\sum_{\substack{0\le j\le S\\j\equiv s\ (r)}}f_j,
 \quad 0\le s<r.
\]
The $R_s$ are the total coefficient masses in the residue classes modulo
$r$. Put $\delta_s=R_s-R_{s-1}$, with subscripts interpreted modulo $r$.
As a formal power series, define
\begin{equation}\label{eq:u-general}
 U(q)=\frac{F(q)}{1-q^r}=\sum_{n\ge0}u_nq^n,
 \qquad u_n=\sum_{\ell\ge0}f_{n-r\ell}.
\end{equation}
As usual, coefficients outside their support are zero.

\begin{lemma}[The periodic tail]\label{lem:tail}
If $n\ge\max\{1,S-r+2\}$, then
$u_n-u_{n-1}=\delta_{n\bmod r}$.
If also $n<rb$, the same difference occurs in
$P=F\qir b r$:
\begin{equation}\label{eq:taildiff}
 \coef nP-\coef{n-1}P=\delta_{n\bmod r}.
\end{equation}
\end{lemma}
\begin{proof}
For the stated $n$, the next possible omitted indices in the residue
classes of $n$ and $n-1$ are $n+r>S$ and $n-1+r>S$.
Consequently $u_n=R_{n\bmod r}$ and $u_{n-1}=R_{(n-1)\bmod r}$.
Finally,
\[
 P(q)=(1-q^{rb})U(q),
\]
so its coefficients agree with those of $U$ below exponent $rb$.
\end{proof}

\begin{theorem}[A necessary degree bound for every spacing]\label{thm:degreebound}
Suppose no $a_i$ is divisible by $r$. If $P_{\boldsymbol a,b,r}$ is
unimodal, then
\begin{equation}\label{eq:necessarydegree}
 r(b-1)\le S,
 \qquad\text{equivalently}\qquad
 b\le1+\left\lfloor\frac Sr\right\rfloor.
\end{equation}
\end{theorem}
\begin{proof}
Let $\zeta$ be a primitive $r$th root of unity.
Every factor
\[
 \qi{a_i}\big|_{q=\zeta}=
       \frac{1-\zeta^{a_i}}{1-\zeta}
\]
is nonzero, so $F(\zeta)\ne0$. The $R_s$ cannot all be equal, since
otherwise $F(\zeta)=R_0(1+\zeta+\cdots+\zeta^{r-1})=0$.
The differences $\delta_s$ sum to zero and are not all zero, so at least
one is negative.

Choose an integer $n$ in the $r$-term interval
\[
 S-r+2\le n\le S+1
\]
whose residue has a negative difference. Such an $n$ is automatically
positive. Indeed, if the interval includes nonpositive integers then
$S\le r-2$; at $n=0$ the difference is $R_0-R_{r-1}=f_0>0$,
and for any negative $n$ in this interval both relevant residue masses
vanish.
Thus Lemma~\ref{lem:tail} applies.

Suppose for a contradiction that $L=r(b-1)>S$.
Then $n\le S+1<rb$. If $L\ge S+2$, we also have
$n\le S+1\le\lfloor(S+L)/2\rfloor$. Equation~\eqref{eq:taildiff}
gives a negative first-half difference, contradicting unimodality.

The remaining possibility is $L=S+1$. In that case
$S\equiv-1\pmod r$. Symmetry of $F$ implies $R_0=R_{r-1}$, hence
$\delta_0=0$. The selected negative difference cannot be at $n=S+1$,
whose residue is zero. Therefore $n\le S=\lfloor(S+L)/2\rfloor$,
and again \eqref{eq:taildiff} contradicts unimodality.
\end{proof}

Together, Theorems~\ref{thm:sufficient} and \ref{thm:degreebound} give a
useful sandwich. When there is no divisible factor,
\begin{equation}\label{eq:sandwich}
 b\le1+\sum_i\left\lfloor\frac{a_i}{r}\right\rfloor
 \ \Longrightarrow\ P\text{ unimodal}\
 \ \Longrightarrow\
 b\le1+\left\lfloor\frac{\sum_i(a_i-1)}r\right\rfloor.
\end{equation}
The residue classes decide what happens in the possible gap.

\section{Spacing two: the original criterion is exact}
\label{sec:r2}
\begin{theorem}[Complete spacing-two criterion]\label{thm:r2}
The polynomial $(\prod_i\qi{a_i})\qir b2$ is unimodal if and only if some
$a_i$ is even or
\[
 b\le1+\sum_i\left\lfloor\frac{a_i}{2}\right\rfloor.
\]
\end{theorem}
\begin{proof}
Sufficiency follows from Theorem~\ref{thm:sufficient}.
If all $a_i$ are odd, write $a_i=2A_i+1$ and $Q=\sum_iA_i$.
Then $S=2Q$, so Theorem~\ref{thm:degreebound} gives the same necessary
bound $b\le Q+1$.
\end{proof}

There is also an exact failure certificate. With all $a_i$ odd,
$F(-1)=1$, so the even residue mass exceeds the odd residue mass by one.
For $b\ge Q+2$, the index $n=S+1=2Q+1$ lies below $2b$ and no later than
the center $Q+b-1$. Lemma~\ref{lem:tail} gives
\[
 \coef{2Q+1}P-\coef{2Q}P=-1.
\]
This argument also covers factors $a_i=1$.

\section{Spacing three: a complete replacement theorem}
\label{sec:r3}
Assume initially that no $a_i$ is divisible by three. Define
\begin{equation}\label{eq:Qt}
 a_i=3A_i+s_i,\quad s_i\in\{1,2\},\qquad
 Q=\sum_iA_i,\qquad t=\#\{i:s_i=2\}.
\end{equation}
Then
\begin{equation}\label{eq:SQt}
 S=3Q+t.
\end{equation}

\begin{theorem}[Sharp spacing-three criterion]\label{thm:r3}
The polynomial $P_{\boldsymbol a,b,3}$ is unimodal if and only if
\begin{equation}\label{eq:r3criterion}
 (\exists i\colon3\mid a_i)
 \quad\text{or}\quad
 b\le1+Q+2\left\lfloor\frac t6\right\rfloor,
\end{equation}
where $Q,t$ are as in \eqref{eq:Qt} in the absence of a divisible factor.
If no $a_i$ is divisible by three and the inequality fails, write
$t=6h+v$, $0\le v\le5$, and $M=Q+2h$. At the explicit index
\begin{equation}\label{eq:dropindex}
 n_0=3M+\left\lfloor\frac{v+3}{2}\right\rfloor
\end{equation}
we have
\begin{equation}\label{eq:unitdrop}
 n_0\le\lfloor D/2\rfloor,\qquad
 \coef{n_0}P-\coef{n_0-1}P=-1.
\end{equation}
\end{theorem}

The proof has two independent parts. A binomial-prefix lemma proves
sufficiency. A six-case residue table proves necessity, including the
unit-drop certificate.

\subsection{The binomial-prefix lemma}
\begin{lemma}[Enough ordinary binary factors]\label{lem:binomialcomb}
For integers $B\ge1$ and $t\ge0$, the polynomial
\[
 (1+q)^t\qir B3
\]
is unimodal whenever $t\ge6\lfloor B/2\rfloor$.
\end{lemma}
\begin{proof}
For $B=1$ this is an ordinary binomial polynomial. Suppose $B\ge2$,
put $m=\lfloor B/2\rfloor$ and $K=6m$, and first take $t=K$.
Thus $B=2m$ or $B=2m+1$.
Let
\[
 f_j=\binom Kj,\qquad
 u_n=\sum_{\ell\ge0}f_{n-3\ell},\qquad d_n=u_n-u_{n-1},
\]
with $f_j=0$ outside $0\le j\le K$ and $u_n=d_n=0$ for $n<0$.
Let $R_0,R_1,R_2$ be the full residue masses of $(1+q)^K$.

Modulo $\Phi_3(q)=1+q+q^2$, one has $(1+q)^3\equiv-1$ and hence
$(1+q)^6\equiv1$. Since $6\mid K$,
\[
 R_0+R_1q+R_2q^2\equiv1\pmod{\Phi_3(q)}.
\]
Two polynomials of degree at most two congruent modulo $\Phi_3$ differ
by a constant multiple of $\Phi_3$. It follows that $R_0-R_1=1$ and
$R_1=R_2$. Since $R_0+R_1+R_2=2^K$,
\begin{equation}\label{eq:binomialresidues}
 R_0=\frac{2^K+2}{3},\qquad
 R_1=R_2=\frac{2^K-1}{3}.
\end{equation}
Consequently
\begin{equation}\label{eq:deltasK}
 (\delta_0,\delta_1,\delta_2)=(1,-1,0).
\end{equation}

Symmetry $f_j=f_{K-j}$ gives a useful reflection identity.
The terms of residue $n$ omitted from $u_n$ have indices
$n+3,n+6,\ldots$. Reflecting these indices about $K/2$ gives
\[
 u_n=R_{n\bmod3}-u_{K-n-3}\qquad(n\ge-1).
\]
Subtract the same identity at $n-1$. For $n\ge0$ this yields
\begin{equation}\label{eq:reflection}
 d_n=\delta_{n\bmod3}+d_{K-n-2}.
\end{equation}

For $0\le n\le K/2$,
\[
 d_n=\sum_{\ell\ge0}\bigl(f_{n-3\ell}-f_{n-1-3\ell}\bigr)\ge1.
\]
Every summand is nonnegative because the binomial coefficients increase
up to $K/2$: explicitly, $f_j/f_{j-1}=(K-j+1)/j>1$ for
$1\le j\le K/2$. The summand with $\ell=0$ is a positive integer.
For $K/2<n\le K-2$, the reflected index $K-n-2$ lies between zero and
$K/2$. Equations~\eqref{eq:deltasK}--\eqref{eq:reflection} therefore give
$d_n\ge-1+1=0$.
At the last two indices,
\[
 d_{K-1}=\delta_2=0,\qquad d_K=\delta_0=1.
\]
We have proved
\begin{equation}\label{eq:prefixmonotone}
 d_n\ge0\qquad(0\le n\le K).
\end{equation}

Now
\[
 (1+q)^K\qir B3=(1-q^{3B})\sum_{n\ge0}u_nq^n.
\]
If $B=2m$, its degree is $2K-3$, its left-middle index is $K-2$,
and $3B=K$. If $B=2m+1$, its degree is $2K$, its middle index is $K$,
and $3B=K+3$. In either case the coefficients through the center agree
with $u_n$, and their differences are nonnegative by
\eqref{eq:prefixmonotone}. Symmetry proves unimodality for $t=K$.
Finally, multiplication by $(1+q)^{t-K}$ preserves $\SU$
(Lemma~\ref{lem:smoothing}), proving the claim for every $t\ge K$.
\end{proof}

\subsection{Sufficiency of the corrected criterion}
\begin{proof}[Proof of Theorem~\ref{thm:r3}, sufficient direction]
A divisible factor is handled by Lemma~\ref{lem:divisible}.
Otherwise assume the right-hand inequality in \eqref{eq:r3criterion}.
For $b\le Q+1$, Theorem~\ref{thm:sufficient} applies.
For $b>Q+1$, use Corollary~\ref{cor:core} with $r=3$.
The residue product is
\[
 (1+q)^t\qir{b-Q}{3}.
\]
Set $B=b-Q$. The assumed inequality says
$B\le1+2\lfloor t/6\rfloor$, which implies
$t\ge6\lfloor B/2\rfloor$.
Lemma~\ref{lem:binomialcomb} proves that the residue product is $\SU$,
and Corollary~\ref{cor:core} lifts this back to the original product.
\end{proof}

\subsection{Necessity and the coefficient-drop certificate}
\begin{proof}[Proof of Theorem~\ref{thm:r3}, necessary direction]
Suppose no $a_i$ is divisible by three, and write
$t=6h+v$, $0\le v\le5$, and $M=Q+2h$.
Then $S=3M+v$ by \eqref{eq:SQt}.

Let $R_0,R_1,R_2$ be the residue masses of
$F=\prod_i\qi{a_i}$ modulo three. Complete blocks of three terms vanish
modulo $\Phi_3(q)=1+q+q^2$, so
\begin{equation}\label{eq:corecongruence}
 F(q)\equiv\prod_i\qi{s_i}=(1+q)^t
       \equiv(1+q)^v\pmod{\Phi_3(q)}.
\end{equation}
The degree-at-most-two residue polynomials of $F$ and $(1+q)^v$
are congruent modulo $\Phi_3$, so they differ by a constant multiple
of $\Phi_3$. Their corresponding residue differences are therefore equal.
The binomial theorem for $v=0,\ldots,5$ gives
Table~\ref{tab:residues}.

\begin{table}[htbp]
\centering
\begin{tabular}{@{}rrrrrrr@{}}
\toprule
$v$ & $\delta_0$ & $\delta_1$ & $\delta_2$
 & $e=\lfloor(v+3)/2\rfloor$ & $n_0\bmod3$ & $n_0-S$ \\
\midrule
0 &  1 & -1 &  0 & 1 & 1 &  1 \\
1 &  1 &  0 & -1 & 2 & 2 &  1 \\
2 &  0 &  1 & -1 & 2 & 2 &  0 \\
3 & -1 &  1 &  0 & 3 & 0 &  0 \\
4 & -1 &  0 &  1 & 3 & 0 & -1 \\
5 &  0 & -1 &  1 & 4 & 1 & -1 \\
\bottomrule
\end{tabular}
\caption{The residue class selected by $n_0=3M+e$ always has difference $-1$.}
\label{tab:residues}
\end{table}

If the bound in \eqref{eq:r3criterion} fails, then $b\ge M+2$.
Choose $n_0$ from \eqref{eq:dropindex}.
The table shows $n_0\ge S-1$ and $n_0\ge1$.
Moreover,
\[
 n_0\le3M+4<3M+6\le3b.
\]
Thus Lemma~\ref{lem:tail} applies at $n_0$, and the selected table entry
gives the exact difference $-1$.
Finally,
\begin{align*}
 \left\lfloor\frac{D}{2}\right\rfloor
 &=\left\lfloor\frac{3M+v+3(b-1)}{2}\right\rfloor\\
 &\ge\left\lfloor\frac{6M+v+3}{2}\right\rfloor
 =3M+\left\lfloor\frac{v+3}{2}\right\rfloor=n_0.
\end{align*}
The negative difference occurs in the first half of the symmetric
sequence, contradicting \eqref{eq:firsthalf}.
This proves both necessity and \eqref{eq:unitdrop}.
\end{proof}

\subsection{Consequences, minimality, and the missing sixth factor}
\label{sec:r3-consequences}
\begin{corollary}[A two-parameter infinite family]\label{cor:binomialexact}
For $t\ge0$ and $B\ge1$,
\begin{equation}\label{eq:binomialexact}
 (1+q)^t\qir B3\text{ is unimodal}
 \quad\Longleftrightarrow\quad
 t\ge6\left\lfloor\frac B2\right\rfloor.
\end{equation}
\end{corollary}
\begin{proof}
For $t\ge1$, take $t$ ordinary factors $a_i=2$ in
Theorem~\ref{thm:r3}, so $Q=0$. The condition becomes
$B\le1+2\lfloor t/6\rfloor$, equivalent to \eqref{eq:binomialexact}.
For $t=0$, the spaced polynomial is unimodal only for $B=1$;
for $B\ge2$ its first coefficients include $1,0$ and a later $1$.
\end{proof}

Every pair $B\ge2$, $t\ge6\lfloor B/2\rfloor$ supplies a counterexample
to the original necessity assertion: the source condition would require
$B\le1$. Thus the degree-nine example is not an isolated coincidence, and
the family is genuinely two-parameter.
Corollary~\ref{cor:infinite} below strengthens the one-parameter slice
$B=2$ to log-concavity; whether log-concavity also holds throughout the
wider two-parameter family of \eqref{eq:binomialexact} is not settled
here and would be a natural next check. The two families are kept side by
side; see item~D4 of Table~\ref{tab:dual}.

\begin{corollary}[Minimality in the spacing-three branch]\label{cor:minimal}
Among counterexamples to the original necessity assertion with $r=3$,
at least six nontrivial ordinary factors are required, and the total
polynomial degree is at least nine. Both bounds are attained by
$(1+q)^6(1+q^3)$.
\end{corollary}
\begin{proof}
A counterexample has no divisible factor and satisfies
\[
 Q+1<b\le Q+1+2\lfloor t/6\rfloor.
\]
Hence $t\ge6$, so there are at least six factors of residue two, each
nontrivial. Also $S=3Q+t\ge6$ and $b\ge2$, whence
$D=S+3(b-1)\ge9$. The displayed example realizes equality.
\end{proof}

This is a statement about the $r=3$ branch only; it asserts nothing about
other spacings. Proposition~\ref{prop:minimum} closes exactly that gap by
a different method, and the two results are retained together; see item~D3
of Table~\ref{tab:dual}.

For $k\le5$, one necessarily has $t\le5$, so the correction
$2\lfloor t/6\rfloor$ vanishes. The corrected theorem therefore agrees
exactly with the old spacing-three criterion throughout that range.
This explains the relevance of the source's reported cutoff of five
ordinary factors \cite[p.~14]{CIMSY}; it is not a numerical accident.

\begin{table}[htbp]
\centering
\small
\begin{tabular}{@{}lrrrrl@{}}
\toprule
Ordinary factors & $Q$ & $t$ & Old bound & Exact bound & Extra admissible $b$ \\
\midrule
$\qi2^6$  & 0 & 6  & 1 & 3 & $2,3$ \\
$\qi2^{12}$ & 0 & 12 & 1 & 5 & $2,3,4,5$ \\
$\qi5^6$ & 6 & 6 & 7 & 9 & $8,9$ \\
$\qi2^5\qi4$ & 1 & 5 & 2 & 2 & None \\
\bottomrule
\end{tabular}
\caption{The exact bound means the largest $b$ for which the product with $\qir b3$ is unimodal.}
\label{tab:examples}
\end{table}

The mechanism is arithmetic: modulo $\Phi_3$, a residue-two factor is
$1+q$, whose sixth power is one. Each complete group of six such factors
allows two additional terms in the spaced factor. In contrast, for
$r=2$ the only nonzero residue is one, so no analogous correction exists.

\section{The two-shift family: an exact smoothing threshold}
\label{sec:twoshift}
The smallest counterexample lives in the slice $b=2$, $a_1=\cdots=a_n=2$
of the general product \eqref{eq:product}. That slice is small enough to
classify completely for \emph{every} spacing, and the classification
explains the special role of the exponent six at $r=3$ from a second,
entirely independent direction. Write
\begin{equation}\label{eq:family}
 \TS_{n,r}(q)=(1+q)^n(1+q^r)=\sum_{j=0}^{n+r}c_j^{(n,r)}q^j,
 \qquad
 c_j^{(n,r)}=\binom nj+\binom n{j-r},
\end{equation}
with binomial coefficients taken to be zero when the lower index leaves
$\{0,\ldots,n\}$. In the notation of \eqref{eq:product},
$\TS_{n,r}=P_{(2,\ldots,2),2,r}$ with $n$ ordinary factors.

\subsection{A weighted subset model}
An object is a pair $(S,\varepsilon)$ with
$S\subseteq\{1,\ldots,n\}$ and $\varepsilon\in\{0,1\}$, of weight
\[
 w(S,\varepsilon)=|S|+r\varepsilon .
\]
Then $c_j^{(n,r)}$ counts objects of weight $j$. Complementation is a
weight-reversing involution,
\[
 (S,\varepsilon)\longmapsto
   (\{1,\ldots,n\}\setminus S,\,1-\varepsilon),
 \qquad w\longmapsto n+r-w,
\]
so that
\begin{equation}\label{eq:symmetry}
 c_j^{(n,r)}=c_{n+r-j}^{(n,r)} .
\end{equation}
This is a weighted enumeration; it is not an assertion that $w$ is the
ordinary rank function on a product of two familiar graded posets.

\subsection{The threshold}
\begin{theorem}[Sharp two-shift smoothing threshold]\label{thm:twoshift}
For integers $r\ge2$ and $n\ge0$,
\begin{equation}\label{eq:twoshift}
 \TS_{n,r}(q)\text{ is unimodal}
 \quad\Longleftrightarrow\quad n\ge r^2-3.
\end{equation}
\end{theorem}

By Lemma~\ref{lem:smoothing} the set of exponents $n$ for which
$\TS_{n,r}$ is unimodal is upward closed, so it suffices to prove failure
at $n=r^2-4$ and unimodality at $n=r^2-3$, each by an argument valid for
every $r$ rather than by a numerical test.
For $r=2$ the two boundary polynomials are $1+q^2$ and $1+q+q^2+q^3$, so
the theorem holds; assume $r\ge3$ from here on.
Sections~\ref{sec:twoshift-lower} and \ref{sec:twoshift-upper} supply the
two halves, and Section~\ref{sec:twoshift-alt} gives a second, independent
proof of the harder half.
The theorem is proved once; everything labelled ``alternative'' or
``cross-check'' below is a redundant confirmation, retained deliberately
(items D1 and D2 of Table~\ref{tab:dual}).

\subsection{The lower bound: a central dip one step below the threshold}
\label{sec:twoshift-lower}
A symmetric unimodal sequence cannot have a central entry smaller than its
neighbour. For $\TS_{n,r}$ that obstruction factors, and the factorization
is valid for every admissible $n$, not only at the boundary exponent.

\begin{lemma}[Central difference in even total degree]\label{lem:central}
Suppose $n\ge r$ and $n+r=2m$. Then
\begin{equation}\label{eq:central-factor}
 c_m^{(n,r)}-c_{m-1}^{(n,r)}
  =\binom nm\,
     \frac{n+2-r^2}{(m+1)(n-m+1)} .
\end{equation}
\end{lemma}
\begin{proof}
By binomial symmetry, $c_m=2\binom nm$ and
$c_{m-1}=\binom n{m-1}+\binom n{m+1}$.
Divide by the positive quantity $\binom nm$ and use consecutive binomial
ratios:
\[
 \frac{c_m-c_{m-1}}{\binom nm}
  =2-\frac{m}{n-m+1}-\frac{n-m}{m+1}
  =\frac{n+2-r^2}{(m+1)(n-m+1)},
\]
the last step using $2m=n+r$.
\end{proof}

The single factor $n+2-r^2$ carries the whole sign, which is where the
quadratic scale $r^2$ comes from algebraically. Now put $n_-=r^2-4$.
For $r\ge3$ we have $n_-\ge r$, and $n_-+r$ is even, so
\eqref{eq:central-factor} applies and gives
\begin{equation}\label{eq:strict-central-dip}
 c_m^{(n_-,r)}-c_{m-1}^{(n_-,r)}
 =-\frac{2\binom{r^2-4}{m}}{(m+1)(r^2-3-m)}<0,
 \qquad m=\frac{r^2+r-4}{2}.
\end{equation}
Hence $\TS_{r^2-4,r}$ is not unimodal. If $\TS_{n,r}$ were unimodal for
some $n<r^2-4$, repeated application of Lemma~\ref{lem:smoothing} would
make $\TS_{r^2-4,r}$ unimodal, a contradiction. This proves the whole
necessary bound $n\ge r^2-3$, including the small values of $n$ at which
the two shifted binomial rows have gaps or disjoint supports.

That last deduction avoids a real edge-case error: applying
\eqref{eq:central-factor} where $\binom nm=0$ would not prove a strict
central dip. The factorization is therefore used only where its positivity
hypotheses hold, and the obstruction is propagated downward by
contraposition of a valid smoothing lemma.

\begin{remark}[Independent boundary cross-check, D2]\label{rem:boundary-crosscheck}
The boundary case of \eqref{eq:strict-central-dip} can also be obtained
directly, without the general lemma. Put $N=n_-+1=r^2-3$ and
$u=(N-r-1)/2$. Writing the central difference through the first-difference
function \eqref{eq:Delta} below and reflecting, it equals
$\Delta_u-\Delta_{u+1}$, and $\binom N{u+1}=\binom Nu(N-u)/(u+1)$ gives
\begin{equation}\label{eq:centralnegative}
 \Delta_u-\Delta_{u+1}
 =\frac{\binom Nu}{N}
   \left((r+1)-(r-1)\frac{N+r+1}{N-r+1}\right)
 =-\frac{4\binom Nu}{N(N-r+1)}<0,
\end{equation}
the denominator being positive for $r\ge3$. One checks
$u=(r^2-r-4)/2=m-r$, and the two expressions agree: at $r=3$ both equal
$-1$, at $r=4$ both equal $-22$, at $r=5$ both equal $-3230$, matching
\path{data/thresholds.csv}. The general lemma is the more useful
statement; the one-line boundary computation is an independent
confirmation of the same sign.
\end{remark}

\subsection{The upper bound: a monotone ratio certificate}
\label{sec:twoshift-upper}
It remains to prove unimodality at $n=r^2-3$. The difficulty is that near
the midpoint of the whole polynomial one binomial row is already decreasing
while its translate is still increasing, so their first differences must be
compared quantitatively.

Fix $n$, put $N=n+1$, and define for every integer $j$
\begin{equation}\label{eq:Delta}
 \Delta_j=\binom nj-\binom n{j-1}
        =\frac{N-2j}{N}\binom Nj
        \qquad(0\le j\le N).
\end{equation}
Consequently
\begin{equation}\label{eq:d-sign-reflection}
 \Delta_j>0\quad\left(0\le j<\tfrac N2\right),
 \qquad \Delta_{N-j}=-\Delta_j,
\end{equation}
and, whenever both sides are positive, consecutive ratios are
\begin{equation}\label{eq:h-ratio}
 h(j):=\frac{\Delta_{j+1}}{\Delta_j}
    =\frac{(n+1-j)(n-1-2j)}{(j+1)(n+1-2j)} .
\end{equation}
The coefficient differences of $\TS_{n,r}$ satisfy
\begin{equation}\label{eq:sum-differences}
 c_j^{(n,r)}-c_{j-1}^{(n,r)}=\Delta_j+\Delta_{j-r}.
\end{equation}

The next identity is the main calculation. It is stated with a free $n$;
only its specialization to $n=r^2-3$ is used.

\begin{lemma}[Ratio identity]\label{lem:ratio}
Let $y=n+1-r-\tau$. Where the indicated denominators are nonzero,
\begin{equation}\label{eq:ratio-identity}
 h(\tau)h(y-1)-1
 =-\frac{r(n+1)\bigl(n+2-r^2+(n-r-2\tau)^2\bigr)}
 {(\tau+1)\,y\,(n+1-2\tau)(n+3-2y)} .
\end{equation}
\end{lemma}
\begin{proof}
By \eqref{eq:h-ratio} the left side equals
\[
 \frac{(n+1-\tau)(n-1-2\tau)(n+2-y)(n+1-2y)}
      {(\tau+1)\,y\,(n+1-2\tau)(n+3-2y)}-1 .
\]
Place the two terms over the common denominator and substitute
$y=n+1-r-\tau$. The resulting numerator factors as
$-r(n+1)\bigl(n+2-r^2+(n-r-2\tau)^2\bigr)$, which is the claim. After
clearing denominators this is a polynomial identity; no approximation is
involved.
\end{proof}

\paragraph{Specialization at the threshold.}
Set $n=r^2-3$ and
\begin{equation}\label{eq:T}
 T=\frac{n+1-r}{2}=\frac{r^2-r-2}{2},
\end{equation}
an integer because $r(r-1)$ is even. The degree $n+r$ is odd and the left
midpoint of $\TS_{n,r}$ is
\begin{equation}\label{eq:mu}
 \mu=\left\lfloor\frac{n+r}{2}\right\rfloor=T+r-1 .
\end{equation}
Let $1\le j\le\mu$. If $j\le(n+1)/2$, both terms of
\eqref{eq:sum-differences} are nonnegative, with $\Delta_{j-r}=0$ when
$j<r$; that part of the left half is already controlled. Suppose instead
\begin{equation}\label{eq:hard-range}
 \frac{n+1}{2}<j\le\mu,
\end{equation}
and write
\[
 \tau=j-r,\qquad y=n+1-j=2T-\tau .
\]
Then $0\le\tau<T<y<(n+1)/2$; for $\tau\ge0$ note that
$(n+1)/2=(r^2-2)/2\ge r$ when $r\ge3$.
Hence $\Delta_\tau$ and $\Delta_y$ are both positive and, by
\eqref{eq:d-sign-reflection} and \eqref{eq:sum-differences},
\begin{equation}\label{eq:compare-d}
 c_j-c_{j-1}=\Delta_\tau-\Delta_y .
\end{equation}
Define, throughout this positive range and extended to $\tau=T$,
\[
 \rho_\tau=\frac{\Delta_\tau}{\Delta_{2T-\tau}},
 \qquad\text{so that}\qquad \rho_T=1 .
\]
For $\tau<T$,
\begin{equation}\label{eq:R-step}
 \frac{\rho_{\tau+1}}{\rho_\tau}
  =\frac{\Delta_{\tau+1}}{\Delta_\tau}\,
     \frac{\Delta_{2T-\tau}}{\Delta_{2T-\tau-1}}
  =h(\tau)h(y-1).
\end{equation}
All four factors of the denominator in \eqref{eq:ratio-identity} are
positive here, and at $n=r^2-3$ its parenthesized numerator becomes
\begin{equation}\label{eq:nonnegative-factor}
 n+2-r^2+(n-r-2\tau)^2
 =-1+\bigl(2(T-\tau)-1\bigr)^2
 =4(T-\tau)(T-\tau-1)\ \ge 0
 \qquad(\tau\le T-1).
\end{equation}
Therefore $\rho_{\tau+1}\le\rho_\tau$, and iterating towards the endpoint
gives $\rho_\tau\ge\rho_T=1$. By \eqref{eq:compare-d} every remaining
left-half difference is nonnegative, so $\TS_{r^2-3,r}$ is unimodal.
Together with Section~\ref{sec:twoshift-lower} and
Lemma~\ref{lem:smoothing}, this proves Theorem~\ref{thm:twoshift}.
Note that the argument controls \emph{every} left-half difference;
positivity of a central difference alone would not have sufficed.

\subsection{An alternative proof of the upper bound}
\label{sec:twoshift-alt}
The following is a second, self-contained proof of unimodality at
$n=r^2-3$. It compares the same two first differences, but runs the
induction in the opposite direction, from a different anchor, and produces
a different factorization. Neither proof is a transcription of the other,
and each is complete on its own; keeping both is deliberate (item~D1 of
Table~\ref{tab:dual}).

\begin{proof}[Second proof of unimodality at $n=r^2-3$]
Here $N=n+1=r^2-2$ and the degree $n+r$ is odd. For a first-half index
$j\le N/2$ both terms of \eqref{eq:sum-differences} are nonnegative. For
the remaining first-half indices put
\[
 s=N+r-2j,\qquad
 u=j-r=\frac{N-r-s}{2},\qquad
 v=N-j=\frac{N-r+s}{2} .
\]
Because $j>N/2$ and $j\le(n+r-1)/2$, the integer $s$ is even and satisfies
$2\le s<r$. Furthermore $0\le u<v<N/2$, so $\Delta_u,\Delta_v>0$, and by
reflection \eqref{eq:sum-differences} equals $\Delta_u-\Delta_v$.
It therefore suffices to prove
\begin{equation}\label{eq:ratio}
 \varrho_s:=\frac{\Delta_u}{\Delta_v}
 =\frac{r+s}{r-s}\cdot
   \frac{\binom N{(N-r-s)/2}}{\binom N{(N-r+s)/2}}
 \ \ge1 .
\end{equation}

At $s=2$,
\[
 \varrho_2=\frac{(r+2)(N-r)(N-r+2)}
                {(r-2)(N+r)(N+r+2)}=1,
\]
because the numerator minus the denominator is
\begin{equation}\label{eq:ratio-base}
 4N(N-r^2+2)=0 .
\end{equation}
For an even $s$ with $s+2<r$, the adjacent ratio is
\begin{equation}\label{eq:ratio-step}
 \frac{\varrho_{s+2}}{\varrho_s}
 =\frac{(r+s+2)(r-s)(N-r-s)(N-r+s+2)}
        {(r+s)(r-s-2)(N+r+s+2)(N+r-s)} .
\end{equation}
All factors in the denominator are positive. Expanding the numerator minus
the denominator gives
\begin{equation}\label{eq:ratio-positive}
 4Nr(N-r^2+s^2+2s+2)=4Nr\,s(s+2)>0,
\end{equation}
where the equality uses $N=r^2-2$. Starting from $\varrho_2=1$, induction
over the even values of $s$ proves \eqref{eq:ratio}. All first-half
differences are nonnegative, and symmetry gives unimodality.
\end{proof}

\begin{remark}[How the two routes differ]
Under $s=2(T-\tau)$ the two comparisons are the same ratio, but the
inductions are genuine mirror images with different content. The proof of
Section~\ref{sec:twoshift-upper} anchors at $\tau=T$, where
$\rho_T=\Delta_T/\Delta_T=1$ holds tautologically, and steps downward in
$s$ using the free-$n$ identity \eqref{eq:ratio-identity} specialized by
\eqref{eq:nonnegative-factor}. The proof above anchors one step further in,
at $s=2$, where equality is not a tautology but the nontrivial identity
\eqref{eq:ratio-base}, and steps upward in $s$ using
\eqref{eq:ratio-positive}. Only the first route states its identity for
arbitrary $n$, which is what makes the equality analysis of
Section~\ref{sec:twoshift-equality} fall out of it; only the second avoids
introducing the auxiliary variable $y$. Their algebraic certificates,
\eqref{eq:ratio-base}--\eqref{eq:ratio-positive} versus
\eqref{eq:ratio-identity}--\eqref{eq:nonnegative-factor}, are checked by
two different computer algebra systems in
Section~\ref{sec:verification-cas}.
\end{remark}

\subsection{Equality, plateaus, and log-concavity}
\label{sec:twoshift-equality}
The equality case is visible in the first proof, with no extra asymptotics
or coefficient calculations.

\begin{corollary}[Exactly four maxima at the threshold]\label{cor:plateau}
Let $r\ge2$, $n=r^2-3$ and $\mu=(n+r-1)/2$. The coefficients of
$\TS_{n,r}$ attain their maximum exactly at
\begin{equation}\label{eq:four-modes}
 \mu-1,\ \mu,\ \mu+1,\ \mu+2,
\end{equation}
and are strictly increasing before this block and strictly decreasing
after it. Empty flanks, as occur for $r=2$, are allowed.
\end{corollary}
\begin{proof}
For $r=2$ the polynomial is $1+q+q^2+q^3$. Assume $r\ge3$.
In \eqref{eq:nonnegative-factor} equality holds at $\tau=T-1$ and the
factor is strictly positive at every earlier integer of the comparison
range. Hence $\rho_{T-1}=\rho_T=1$ while $\rho_\tau>1$ for $\tau<T-1$, so
the only vanishing left-half difference is at $j=(T-1)+r=\mu$.
In the easy range $j\le(n+1)/2$ the first term $\Delta_j$ is strictly
positive unless $j=(n+1)/2$ is an integer; at that exceptional index
$\Delta_{j-r}>0$, so the sum is still strictly positive.
Consequently $c_{\mu-1}=c_\mu$ and all preceding differences are positive.
Odd-degree symmetry gives $c_\mu=c_{\mu+1}$ and $c_{\mu+2}=c_{\mu-1}$.
\end{proof}

The plateau then contracts under further smoothing. By
\eqref{eq:smoothing-difference}, a flat maximum block of length $w\ge2$
with strict flanks becomes a block of length $w-1$, and a unique maximum
becomes a two-term central block. Starting at $n=r^2-3$, the numbers of
maximizing coefficients are therefore
\begin{equation}\label{eq:plateau-evolution}
 4,\ 3,\ 2,\ 1,\ 2,\ 1,\ 2,\ldots
\end{equation}
as $n$ increases by one; the endpoints of each block are fixed by symmetry
and the total degree.

For the next statement, call a nonnegative finite sequence
\emph{log-concave} when it has no internal zeros and $a_j^2\ge a_{j-1}a_{j+1}$
at every internal index. The support condition matters: local inequalities
alone can hold vacuously across a sufficiently long gap.

\begin{lemma}[$(1+q)$ preserves log-concavity]\label{lem:lc}
Multiplication by $1+q$ preserves log-concavity of a positive coefficient
sequence.
\end{lemma}
\begin{proof}
Extend $a_j$ by zero outside its support and set $b_j=a_j+a_{j-1}$. Direct
expansion gives
\begin{equation}\label{eq:lc-identity}
 b_j^2-b_{j-1}b_{j+1}
  =(a_j^2-a_{j-1}a_{j+1})
   +(a_{j-1}^2-a_{j-2}a_j)
   +(a_{j-1}a_j-a_{j-2}a_{j+1}).
\end{equation}
The first two summands are nonnegative by log-concavity. In the interior,
log-concavity says that the positive consecutive ratios $a_j/a_{j-1}$ are
nonincreasing; in particular $a_{j-1}/a_{j-2}\ge a_{j+1}/a_j$, which makes
the third summand nonnegative. At the endpoints the same claim follows
directly from the zero extension, and the new support has no gaps.
\end{proof}

This is a special case of familiar convolution closure phenomena; see
\cite{Stanley} for broader context. The elementary identity above supplies
all that is needed here.

\begin{corollary}[A log-concave infinite family]\label{cor:infinite}
For every $n\ge6$, the polynomial
$\qi2^{\,n}\qir23=(1+q)^n(1+q^3)$ is a log-concave counterexample to the
necessity assertion of \cite[Conjecture 5.4]{CIMSY}.
\end{corollary}
\begin{proof}
At $n=6$ the eight interior quantities $c_j^2-c_{j-1}c_{j+1}$ are
\begin{equation}\label{eq:lc-witness}
 21,\ 99,\ 126,\ 0,\ 0,\ 126,\ 99,\ 21,
\end{equation}
so the witness is log-concave, and Lemma~\ref{lem:lc} propagates this to
every larger $n$. Theorem~\ref{thm:twoshift} supplies unimodality.
For every member $r=3$, no $a_i=2$ is divisible by $r$, and the floor sum
in \eqref{eq:condition} is zero, so the required inequality would read
$2\le1$.
\end{proof}

Corollaries~\ref{cor:binomialexact} and \ref{cor:infinite} are kept
together: the first is broader in parameters, the second stronger for each
member (item D4 of Table~\ref{tab:dual}).
For $r=3$, Theorem~\ref{thm:twoshift} gives the threshold $n=6$, in
agreement with Corollary~\ref{cor:binomialexact} at $B=2$: two independent
routes to the same special value six.

\subsection{Threshold values and a scope warning}
The first thresholds are consequences of Theorem~\ref{thm:twoshift}, not
extrapolations from the table.
\begin{center}
\small
\begin{tabular}{rrrr}
\toprule
Gap $r$ & First unimodal exponent $r^2-3$ & Degree & Four mode indices\\
\midrule
2&1&3&0--3\\
3&6&9&3--6\\
4&13&17&7--10\\
5&22&27&12--15\\
6&33&39&18--21\\
7&46&53&25--28\\
8&61&69&33--36\\
9&78&87&42--45\\
10&97&107&52--55\\
\bottomrule
\end{tabular}
\end{center}
For $r\ge4$ these members still violate the floor-sum condition
\eqref{eq:condition}, but they have both $k=n>3$ and $r>3$. They are
therefore \emph{not} additional counterexamples inside the conjecture's
necessity range; they describe the smoothing phenomenon beyond that range.

\section{Minimum degree across all spacings}
\label{sec:minimal}
A small counterexample need not be the smallest.
Corollary~\ref{cor:minimal} proved, by theory and without computation,
that within the $r=3$ branch at least six nontrivial ordinary factors and
degree at least nine are required, and explicitly said nothing about other
spacings. This section closes that gap by a complete finite search over
\emph{all} spacings, without imposing arbitrary bounds on $k$ or on the
factor lengths.

\subsection{Reduction to finitely many normalized parameter sets}
For a counterexample to necessity we may assume $b\ge2$, since $b=1$
always satisfies \eqref{eq:condition}. Factors $\qi1=1$ can be deleted:
deleting them changes neither the polynomial nor $\mathcal C_r$, and cannot
invalidate the applicability of the necessity assertion, because $k\le3$
remains true when $k$ decreases and $r\le3$ is untouched.
At least one nontrivial ordinary factor must remain, since otherwise the
polynomial is $\qir br$, which is not unimodal for $r,b\ge2$. We may
therefore normalize by
\[
 2\le a_1\le a_2\le\cdots\le a_k .
\]
For fixed $D$, every normalized candidate then satisfies
\begin{equation}\label{eq:finite-bounds}
 2\le r\le D-1,\qquad
 2\le b\le1+\floor{\frac{D-1}{r}},
\end{equation}
and once $r,b$ are fixed the positive integers $a_i-1$ form a partition of
$S=D-r(b-1)\ge1$. Conversely each such partition gives exactly one
normalized parameter set. Hence the enumeration is complete and contains
no permutation duplicates.

\subsection{The certificate}
\begin{proposition}[Exact finite computation]\label{prop:minimum}
Among products \eqref{eq:product} in the asserted necessity range, the
minimum degree of a counterexample to necessity is nine. After deleting
trivial factors and permuting the ordinary factors, the degree-nine
counterexample parameter set is unique, namely $r=3$, $b=2$,
$a_1=\cdots=a_6=2$.
\end{proposition}
\begin{proof}[Finite verification and its completeness argument]
The reductions above exhaust all possible counterexamples of smaller
degree. Running two independent coefficient implementations over those
parameter sets gives the counts below.
\begin{center}
\small
\begin{tabular}{rrrr}
\toprule
Degree & Normalized parameter sets & Necessity counterexamples
 & Sufficiency counterexamples\\
\midrule
1&0&0&0\\
2&0&0&0\\
3&1&0&0\\
4&3&0&0\\
5&7&0&0\\
6&13&0&0\\
7&23&0&0\\
8&38&0&0\\
9&59&1&0\\
\midrule
Total&144&1&0\\
\bottomrule
\end{tabular}
\end{center}
The single flagged row has $r=3$, $b=2$ and six factors of length two.
The complete rows, including all coefficient arrays and both logical flags,
are in \path{data/degree_search.csv}; the executable enumeration is
\path{code/verify_threshold.py}. Appendix~\ref{app:minimal-code} gives a
second, compact, standalone realization of the same search.
\end{proof}

The zero column for sufficiency is consistent with
Theorem~\ref{thm:sufficient}, which proves sufficiency outright; it is a
bounded computational observation, not the reason sufficiency holds.
This is a computer-checked finite enumeration, not a proof-assistant
formalization: its finite coverage is proved above, its arithmetic is
integer-exact, and its whole state space is small enough to inspect.
The threshold theorem of Section~\ref{sec:twoshift} is independent of it.

Proposition~\ref{prop:minimum} and Corollary~\ref{cor:minimal} are kept
side by side because neither implies the other (item~D3 of
Table~\ref{tab:dual}). The corollary explains \emph{why} six residue-two
factors are forced at spacing three, by an argument valid without any
computation; the proposition establishes that no other spacing does better,
which no part of the $r=3$ analysis can see.

\section{Why a quadratic threshold is natural: a heuristic}
\label{sec:probability}
\noindent\textbf{This section is explanatory only. Nothing in it is used in
any proof above, and none of it is offered as evidence for
Theorem~\ref{thm:twoshift}.}

Choose a weighted subset object of Section~\ref{sec:twoshift} uniformly at
random. Its weight is $X+rB$ with $X\sim\operatorname{Bin}(n,1/2)$ and
$B\sim\operatorname{Bernoulli}(1/2)$ independent, and its probability
generating function is $2^{-(n+1)}\TS_{n,r}(q)$. Conditionally on $B$ it is
one of two translated binomial distributions, with means separated by $r$
and with component standard deviation $\sqrt n/2$.

There is an instructive continuous analogue. Up to a positive
normalization, a symmetric mixture of two normal densities with variance
$\sigma^2$ and means $\pm a$ has density
\[
 f(x)=e^{-(x-a)^2/(2\sigma^2)}+e^{-(x+a)^2/(2\sigma^2)}
     =2e^{-(x^2+a^2)/(2\sigma^2)}\cosh\!\left(\frac{ax}{\sigma^2}\right),
\]
so that for $x>0$
\[
 \frac{f'(x)}{f(x)}=\frac{-x+a\tanh(ax/\sigma^2)}{\sigma^2}.
\]
If $a\le\sigma$, then $\tanh u\le u$ for $u\ge0$ gives
$a\tanh(ax/\sigma^2)\le a^2x/\sigma^2\le x$, so the density is
nonincreasing on the positive half-line and is unimodal. If $a>\sigma$,
differentiation gives $f''(0)/f(0)=(a^2-\sigma^2)/\sigma^4>0$, so the
centre is a local minimum rather than a mode. The continuous threshold is
therefore exactly $a=\sigma$. Putting $a=r/2$ and $\sigma=\sqrt n/2$
predicts the scale $n\approx r^2$, against the exact discrete answer
$n\ge r^2-3$.

The quadratic scale thus has a simple variance interpretation, while the
integer correction and the four-term equality plateau require the exact
calculation. Approximation by normal densities cannot rule out small local
coefficient dips, and weak convergence of distributions does not preserve
unimodality in the direction a finite coefficient theorem needs.

\section{Exact computations and reproducibility}
\label{sec:verification}
The proofs above establish the infinite statements. Two separate exact
verifiers check finite cases using Python integers, without floating-point
arithmetic or external packages.
\path{code/verify_products.py} covers the general product
\eqref{eq:product}; the run recorded in
\path{data/verification_results.json} passed every check of
Table~\ref{tab:tests}.
\path{code/verify_threshold.py} covers the two-shift family and the
minimum-degree search; the run recorded in
\path{data/verification_report.json} passed every check of
Table~\ref{tab:tests2}. The two suites have disjoint test families and
disjoint programming interfaces, and neither imports the other.

\begin{table}[htbp]
\centering
\small
\begin{tabular}{@{}p{.69\linewidth}r@{}}
\toprule
Check in \path{code/verify_products.py} & Number of cases \\
\midrule
Independent schoolbook versus optimized multiplication & 1,000 \\
Centered peeling identity on an integer grid & 3,003 \\
Exhaustive spacing-three multiset products & 30,539 \\
Deterministic random spacing-three products, no divisible factor & 30,000 \\
Spacing-three products with a divisible factor & 2,000 \\
General-spacing sufficient condition & 10,000 \\
General-spacing necessary degree bound & 10,000 \\
Spacing-two necessary-and-sufficient criterion & 5,000 \\
Two-shift boundary polynomials, $2\le r\le60$ & 177 \\
Exact two-shift negative central-difference formula & 58 \\
Spacing-three binomial-comb threshold samples & 237 \\
\bottomrule
\end{tabular}
\caption{Finite checks actually executed by the general-product verifier. Some families overlap; the counts are not
numbers of distinct mathematical theorems or independent proofs.}
\label{tab:tests}
\end{table}

\begin{table}[htbp]
\centering
\small
\begin{tabular}{@{}p{.42\linewidth}r p{.38\linewidth}@{}}
\toprule
Check in \path{code/verify_threshold.py} & Count & Range or purpose\\
\midrule
Full rectangular threshold grid&3,263&$2\le r\le14$, $0\le n\le250$\\
Boundary and nearby rows&594&$2\le r\le100$, $n=r^2-3+\delta$, $\delta=-1,\ldots,4$\\
Central-difference identity&98&One exact obstruction for each $3\le r\le100$\\
Ratio comparisons&2,450&All comparisons used at the tested thresholds\\
Minimum-degree parameter sets&144&Every normalized case with $D\le9$\\
Independent tuple enumerations&144&4,762 objects in total\\
\bottomrule
\end{tabular}
\caption{Finite checks actually executed by the two-shift verifier. Some boundary
cases also occur in the rectangular grid; these are test counts, not a claim that
every test involves a distinct parameter pair.}
\label{tab:tests2}
\end{table}

\subsection{What was checked}
The exhaustive spacing-three test considers all multisets of ordinary
parameters from $\{1,2,4,5,7,8\}$, with zero through eight factors,
and all $b$ from one through the corrected bound plus two.
The zero-factor cases check the harmless empty-product boundary.
In this finite collection, 938 products violate the old necessity
condition while satisfying the corrected theorem.
Every predicted nonunimodal case is also checked at the explicit index
\eqref{eq:dropindex}; the drop is exactly $-1$, not merely negative.

The separate random spacing-three experiment uses seed 1977, chooses
$k$ uniformly from $1,\ldots,16$, chooses each $a_i$ from the integers
$1,\ldots,69$ not divisible by three, and chooses $b$ uniformly from
one through the corrected bound plus eight. The seed is reset for this
experiment so it is independent of how many preliminary tests are run.

The two-shift boundary checks additionally verify the plateau lengths
$4,3,2,1,2$ predicted by \eqref{eq:plateau-evolution} and the strictness
of the threshold's left flank.

\subsection{Five coefficient routes and two unimodality tests}
The small counterexample is expanded by three mutually independent routes
in \path{code/verify_products.py} alone: the sliding-sum recurrence, the
binomial theorem, and direct enumeration of all $2^7=128$ choices of six
weights one and one weight three. Across the package, five distinct
coefficient implementations are retained (item~D7 of
Table~\ref{tab:dual}):
\begin{enumerate}[label=(\roman*)]
\item the sliding-sum recurrence
  \begin{equation}\label{eq:sliding}
   v_j=u_j+v_{j-s}-u_{j-as},
  \end{equation}
  obtained by extracting coefficients from
  $(1-q^s)B(q)=(1-q^{as})A(q)$, where $B=\qir as A$;
\item schoolbook convolution, cross-checked against the optimized version
  in 1,000 cases;
\item direct tuple enumeration over $(x_1,\ldots,x_k,z)$ with
  $0\le x_i<a_i$ and $0\le z<b$, incrementing the coefficient at
  $x_1+\cdots+x_k+rz$, using no polynomial multiplication at all
  (144 cases, 4,762 tuples);
\item the binomial theorem for the two-shift family; and
\item the $128$-element binary-choice enumeration of the specific
  counterexample.
\end{enumerate}
Unimodality itself is tested in two ways throughout, and both are kept
(item~D8): a generic detector that looks for a rise after a descent and
does \emph{not} assume symmetry, and a test that uses symmetry plus the
left-half first differences \eqref{eq:firsthalf}.

For a polynomial of degree $D$, one application of \eqref{eq:sliding}
costs $O(D)$ integer additions and subtractions, so a product with $k$
ordinary factors can be formed in $O(kD)$ arithmetic operations and $O(D)$
working storage. These are arithmetic-operation bounds; the bit cost also
depends on coefficient size. For the two-shift family a binomial row is
generated by $\binom n{j+1}=\binom nj(n-j)/(j+1)$ with a checked exact
division at each step, after which one shifted addition produces
\eqref{eq:family}.

\subsection{Two computer algebra witnesses}
\label{sec:verification-cas}
The two proofs of Theorem~\ref{thm:twoshift} rest on different algebraic
certificates, and each was checked on a different system (item~D6).

A Wolfram Language evaluation, using the connected Wolfram kernel
reporting version \texttt{15.0.1}, confirmed the counterexample
coefficients, the peeling identity \eqref{eq:peeling}, the six-case
residue table of Table~\ref{tab:residues}, the base identity
\eqref{eq:ratio-base} in the factored form $4nn(2+nn-rr^2)$, the step
identity \eqref{eq:ratio-positive} in the factored form
$4\,nn\,rr\,(2+nn-rr^2+2ss+ss^2)$, and the two-shift boundary tests for
$2\le r\le16$. The exact input and returned text are
\path{code/verify_products.wl} and \path{data/wolfram_output.txt}.

A separate optional script, \path{code/symbolic_checks.py}, uses
SymPy~1.14.0 to simplify three expressions to zero: the ratio identity
\eqref{eq:ratio-identity}, the central-difference identity
\eqref{eq:central-factor}, and the boundary factorization
\eqref{eq:nonnegative-factor}. Its output is
\path{data/symbolic_report.json}; the main verifiers do not require
SymPy.

Both are independent computational checks, not proof-assistant
verification of the arguments, and the symbolic check is an audit of
algebra rather than a substitute for establishing the
positive-denominator hypotheses used in the proofs.

\subsection{Two threshold tables on orthogonal axes}
The package carries two threshold tables, neither a subset of the other
(item~D9). \path{data/r3_binomial_thresholds.csv} is indexed by
$b=1,\ldots,40$ and gives, for spacing three, the least unimodal exponent
$t$ of Corollary~\ref{cor:binomialexact} together with the gap index and
the gap value $-1$ at the preceding exponent.
\path{data/thresholds.csv} is indexed by $r=2,\ldots,100$ and gives
\texttt{critical\_n}, \texttt{below\_n}, the exact
\texttt{below\_central\_difference}, \texttt{critical\_degree} and
\texttt{plateau\_start}/\texttt{plateau\_end}. Both are regenerated by
their respective verifiers.

\subsection{Reproduction commands and output}
From this directory, with Python 3.10 or later:
\begin{lstlisting}[language={}]
python3 code/verify_products.py --full --out data/verification_results.json
python3 code/verify_threshold.py
python3 code/minimal_certificate.py
python3 code/symbolic_checks.py   # optional: requires SymPy
\end{lstlisting}
The first command writes a JSON run summary, the small counterexample
certificate, and the spacing-three threshold CSV; it returns a nonzero exit
code if any check fails. For a shorter smoke test, omit \texttt{--full} and
choose a different output filename so that the recorded full run is
preserved. The second command regenerates the JSON and CSV files of the
two-shift analysis in \path{data/}. All checks use explicit exceptions
rather than \texttt{assert}, so they remain active under \texttt{python -O}.
No network access is used by any verification program.

Build this report with the supplied \texttt{Makefile} or run
\begin{lstlisting}[language={}]
latexmk -pdf -interaction=nonstopmode report.tex
\end{lstlisting}
A TeX installation with the packages listed in the preamble is needed, and
\path{code/minimal_certificate.py} must be present because
Appendix~\ref{app:minimal-code} typesets it directly. No custom font files,
proprietary packages, BibTeX run or shell escape are required.

\subsection{Closed-form decision versus coefficient expansion}
Theorem~\ref{thm:r3} gives a decision procedure requiring $O(k)$
integer-arithmetic operations: test divisibility, compute $Q$ and $t$,
and compare $b$ with the bound. If the answer is negative,
\eqref{eq:dropindex} gives the witness index directly.
No expansion of the polynomial is needed to decide the question.
The actual bit complexity depends on the sizes of the input integers.

The positive proof uses induction and the centered identity.
Literally unrolling all $Q$ peeling steps can be long when the parameters
are encoded in binary; no claim of a uniformly short expanded
certificate is being made. The full finite verifier expands coefficients
because it is intended as an independent check, not as the fastest
implementation of the classification theorem.

\subsection{Deliberate duplications}
The merge that produced this note kept a number of results, proofs and
artefacts in duplicate on purpose. Table~\ref{tab:dual} lists them, with
one line each on what the redundancy buys.

\begin{longtable}{@{}p{.055\linewidth}p{.35\linewidth}p{.53\linewidth}@{}}
\caption{Results, proofs and artefacts deliberately kept in more than one copy.}
\label{tab:dual}\\
\toprule
& Duplicated item & What the second copy buys \\
\midrule
\endfirsthead
\multicolumn{3}{@{}l@{}}{\small\itshape Table \thetable{} continued.}\\
\toprule
& Duplicated item & What the second copy buys \\
\midrule
\endhead
\bottomrule
\endlastfoot
D1 & Two proofs of Theorem~\ref{thm:twoshift}:
Section~\ref{sec:twoshift-upper} (free-$n$ ratio identity, induction on
$\tau$ down from $\rho_T=1$) and Section~\ref{sec:twoshift-alt}
(induction on even $s$ up from $\varrho_2=1$).
& The first states its identity for arbitrary $n$ and therefore yields the
equality analysis for free; the second needs no auxiliary variable and
anchors on a nontrivial identity, so a slip in either factorization cannot
propagate to both. \\
D2 & Two accounts of the failure at $n=r^2-4$: general
Lemma~\ref{lem:central} and the boundary computation
\eqref{eq:centralnegative}.
& The lemma exhibits $n+2-r^2$ as the sign-determining factor for every
admissible $n$; the one-line boundary computation confirms the same sign
independently and agrees numerically. \\
D3 & Two minimality results: Corollary~\ref{cor:minimal} (theory, $r=3$
only) and Proposition~\ref{prop:minimum} (exhaustive, all $r$).
& The first explains why six residue-two factors are forced and needs no
computation; the second rules out every other spacing, which the $r=3$
analysis cannot see. Neither implies the other. \\
D4 & Two counterexample families:
Corollary~\ref{cor:binomialexact} ($(1+q)^t\qir B3$, two parameters) and
Corollary~\ref{cor:infinite} ($(1+q)^n(1+q^3)$, log-concave).
& The first is broader in parameters, the second stronger per member. \\
D5 & Two witness polynomials: the plateau witness $(1+q)^6(1+q^3)$ and the
strict witness $(1+q)^9(1+q^3)$.
& The second has a unique maximum and strictly monotone flanks, answering
the objection that the refutation is a plateau accident. \\
D6 & Two computer algebra witnesses: Wolfram Language 15.0.1
(\path{code/verify_products.wl}) and SymPy 1.14.0
(\path{code/symbolic_checks.py}).
& They check the two \emph{different} algebraic routes of D1 on two
different systems. \\
D7 & Five coefficient implementations, listed in
Section~\ref{sec:verification}.
& Agreement between a recurrence, a convolution, a pure counting
enumeration and a closed form makes a shared implementation bug very
unlikely. \\
D8 & Two unimodality tests: generic rise-after-descent, and symmetry plus
left-half differences.
& The generic test does not assume symmetry, so it would catch an error in
the symmetry reasoning itself. \\
D9 & Two threshold tables:
\path{data/r3_binomial_thresholds.csv} (axis $b$) and
\path{data/thresholds.csv} (axis $r$).
& Orthogonal axes; neither table's rows are derivable from the other's. \\
D10 & Two closure lemmas: Lemma~\ref{lem:closure} (centered intervals) and
Lemma~\ref{lem:smoothing} ($b_j-b_{j-1}=a_j-a_{j-2}$).
& The first is the general tool the peeling induction needs, with its
common-centre convention; the second is a short special case sufficient for
Section~\ref{sec:twoshift}, and it also supplies
\eqref{eq:plateau-evolution}. \\
D11 & Two provenance records, merged in \texttt{PROVENANCE.md}.
& The union of search queries, pinpoints, dates and scope caveats from both
original packages; see Section~\ref{sec:provenance}. \\
\end{longtable}

\section{Provenance of this merged note}
\label{sec:provenance}
This note is the merge of two independently produced research packages
that attacked the same conjecture and arrived at the same theorem.

\path{docs/reports/log-concavity-and-unimodality/q-integer-product-unimodality}
(18 September 2026) treated the general product \eqref{eq:product}. It
proved the shared threshold theorem as a supplementary result, through the
first-difference function $H_n(j)=\binom nj-\binom n{j-1}=((N-2j)/N)\binom Nj$
with $N=n+1$, an induction over the even parameter $s=N+r-2j$ in $[2,r)$
on the ratio $H_n(u)/H_n(v)$, anchored at $s=2$ by the identity
$4N(N-r^2+2)=0$ and stepped by $4Nr\,s(s+2)>0$; and it computed the failure
at $n=r^2-4$ as the single central difference $-4\binom Nu/(N(N-r+1))$.
That route survives here as Section~\ref{sec:twoshift-alt} and
Remark~\ref{rem:boundary-crosscheck}. The present note keeps that package's
spine: Sections~\ref{sec:tools}--\ref{sec:r3} and
Sections~\ref{sec:verification}--\ref{sec:scope} descend from it, including
the centered peeling identity, the residue-mass obstruction, the
spacing-two and spacing-three classifications, and the open-problem
programme for $r\ge4$.

\path{docs/reports/log-concavity-and-unimodality/binomial-smoothing-threshold}
(20 September 2026) treated only the two-shift family, and made the shared
threshold its main theorem. It proved it through
$d_j=\binom nj-\binom n{j-1}=((n+1-2j)/(n+1))\binom{n+1}{j}$, the
consecutive ratio $h(j)=d_{j+1}/d_j$, and the free-$n$ ratio identity
\eqref{eq:ratio-identity}, running an induction on $\tau\le T=(r^2-r-2)/2$
over $R_\tau=d_\tau/d_{2T-\tau}$ anchored at $R_T=1$ and using the
specialization $4(T-\tau)(T-\tau-1)\ge0$; and it proved the general
central-difference Lemma~\ref{lem:central} for every $n\ge r$ with $n+r$
even, specializing it at the boundary. That route is the primary proof
here, Sections~\ref{sec:twoshift-lower}--\ref{sec:twoshift-upper}, because
its free-$n$ statement yields Corollary~\ref{cor:plateau} directly.
Sections~\ref{sec:twoshift}, \ref{sec:minimal} and \ref{sec:probability},
and Appendices~\ref{app:minimal-code} and \ref{app:random}, descend from
that package.

Both packages refuted the source's necessity assertion with the identical
witness and identical parameters, and both reduced the threshold to the
two adjacent exponents $n=r^2-3$ and $n=r^2-4$ by the same upward-closure
argument. The shared theorem is stated and proved once here, as
Theorem~\ref{thm:twoshift}; the second route is retained explicitly as an
alternative proof. Search records, pinpoints and scope caveats from both
packages are merged in \texttt{PROVENANCE.md} (item~D11 of
Table~\ref{tab:dual}). Neither package was refereed
or machine-verified, and no exhaustive priority determination was made by
either.

\section{Scope and further targets}
\label{sec:scope}
The exact counterexample disposes of the stated necessity assertion.
The stronger results retain and prove its sufficient half, classify all
spacing-two and spacing-three products, classify the two-shift family for
every spacing, and identify a precise residue obstruction rather than a
merely numerical discrepancy. The proofs explain both the first
counterexample and the complete families that it begins.

The source's main $q$-Fibonomial question is different. This note neither
proves nor disproves that main conjecture, and it does not challenge the
source's established small-case theorems. The counterexample addresses the
auxiliary product criterion \eqref{eq:condition}, not a claim that a
particular $q$-Fibonomial coefficient is nonunimodal.

For $r\ge4$, a general classification of
$(\prod_i\qi{a_i})\qir b r$ is not obtained here. In particular, the
necessity question for two or three ordinary factors at those spacings
is not settled by the present arguments.
Restricting the source's necessity clause to the surviving regimes would
produce a new statement needing its own proof, not an automatic repair.
What can be said precisely is this: necessity is \emph{true} for $r=2$
(Theorem~\ref{thm:r2}) and for $r=3$ whenever $k\le5$
(Theorem~\ref{thm:r3}, since then $t\le5$), \emph{false} for $r=3$ with
$t\ge6$, and \emph{open} here for $k\le3$ with $r\ge4$.

A concrete continuation is to combine the centered reduction with the
full vector of residue differences at larger $r$.
Lemma~\ref{lem:tail} gives exact obstruction indices whenever the
periodic tail enters the first half. Corollary~\ref{cor:core} gives
positive constructions from residue products.
For $r=3$, these two mechanisms meet exactly because
$(1+q)^6\equiv1\pmod{\Phi_3(q)}$.
For larger $r$, the residual factors $\qi{s_i}$ range over more than
$1$ and $1+q$, so the same one-parameter six-case reduction is no longer
available. Determining when the positive and negative tests meet is a
specific next task, not a conclusion asserted in this note.

Two further questions are left open by the merge itself: whether the wider
two-parameter family of Corollary~\ref{cor:binomialexact} is log-concave
throughout, as the slice $B=2$ is by Corollary~\ref{cor:infinite}; and
whether the minimum-degree uniqueness of Proposition~\ref{prop:minimum}
persists at degrees ten and above, where the finite search was not run.

The structural lesson is narrower and more useful than a blanket rejection
of arithmetic criteria: an invariant that assigns zero contribution to all
short factors can fail to detect their cumulative smoothing effect. For
Bernoulli factors that effect is measurable exactly --- a gap of size $r$
disappears at exponent $r^2-3$ --- and the factored ratio identity explains
both the threshold and its unusually flat equality case.

\appendix
\section{A short standalone certificate}
\label{app:short-certificate}
The disproof itself does not depend on any of the stronger theorems.
The following elementary Python fragment checks it using only the
binomial theorem.
\begin{lstlisting}[language=Python]
from math import comb

c = [
    (comb(6, j) if 0 <= j <= 6 else 0)
    + (comb(6, j-3) if 0 <= j-3 <= 6 else 0)
    for j in range(10)
]
assert c == [1, 6, 15, 21, 21, 21, 21, 15, 6, 1]
assert c == c[::-1]
assert all(c[j-1] <= c[j] for j in range(1, 5))
a, b, r = [2]*6, 2, 3
assert not any(x % r == 0 for x in a)
assert b > 1 + sum(x//r for x in a)
assert len(a) <= 3 or r <= 3
\end{lstlisting}
The distributed full verifiers use explicit exceptions instead of
\texttt{assert}, so running Python with optimization does not disable their
checks. The fragment above is only a compact readable certificate.

\section{A compact independent minimum-degree verifier}
\label{app:minimal-code}
The following program uses direct tuple enumeration only. It checks all
144 normalized cases through degree nine, prints their degree counts, and
requires exactly the one stated violation. It has no third-party
dependencies, and its explicit exception checks also remain active under
Python's optimization flag. It is independent of
\path{code/verify_threshold.py}: it shares no code with it and uses no
polynomial multiplication.

\lstinputlisting[language=Python,basicstyle=\ttfamily\footnotesize]{code/minimal_certificate.py}

\section{Recorded random area selection}
\label{app:random}
The second of the two merged packages began from a recorded random draw
over research areas. The initial tool execution was
\begin{lstlisting}[language=Python]
index = secrets.randbelow(len(areas))
selected_area = areas[index]
\end{lstlisting}
It returned the zero-based index $0$, selecting enumerative combinatorics.
There was no redraw. The draw used Python's \texttt{secrets} module rather
than a manually chosen seed. The raw operating-system entropy was not
recorded; the realized index and the exact ordered list are saved in
\path{data/area_selection.json}. Re-running the draw is not expected to
reproduce the same index. All mathematical calculations in this note are
deterministic.

\begin{center}
\small
\begin{tabular}{r l r l}
\toprule
Index&Area&Index&Area\\
\midrule
0&Enumerative combinatorics&12&Algebraic combinatorics\\
1&Graph theory&13&Automata and formal languages\\
2&Combinatorics on words&14&Probability on finite structures\\
3&Finite geometry&15&Topological graph theory\\
4&Number theory&16&Extremal set theory\\
5&Permutation patterns&17&Polyhedral combinatorics\\
6&Order theory&18&Combinatorial game theory\\
7&Semigroup theory&19&Polynomial inequalities\\
8&Discrete geometry&20&Matrix theory\\
9&Additive combinatorics&21&Coding theory\\
10&Symbolic dynamics&22&Partition theory\\
11&Recreational integer sequences&23&Asymptotic analysis\\
\bottomrule
\end{tabular}
\end{center}

\begin{thebibliography}{9}
\bibitem{CIMSY}
Brendan B. Connelly, Ezekiel Ito, Thomas C. Martinez,
Olha Shevchenko, and Kacey Yang.
\newblock \emph{Unimodality of $q$-Fibonomial coefficients for small cases}.
\newblock arXiv:2605.12822v1 [math.CO], submitted 12 May 2026.
\newblock \href{https://arxiv.org/abs/2605.12822v1}{Version-pinned arXiv record};
\href{https://arxiv.org/html/2605.12822v1}{HTML};
\href{https://arxiv.org/pdf/2605.12822v1}{PDF}.
\newblock The target is Section 5.2, Conjecture 5.4 and the adjacent testing
statement, p.~14.

\bibitem{BCK}
Nantel Bergeron, Cesar Ceballos, and Josef K\"ustner.
\newblock Elliptic and $q$-analogs of the Fibonomial numbers.
\newblock \emph{SIGMA} \textbf{16} (2020), paper 076, 16 pp.
\newblock \href{https://doi.org/10.3842/SIGMA.2020.076}{doi:10.3842/SIGMA.2020.076}.
\newblock Background for the distinction from the original Fibonomial question.

\bibitem{Stanley}
Richard P. Stanley.
\newblock Log-concave and unimodal sequences in algebra, combinatorics,
and geometry.
\newblock \emph{Annals of the New York Academy of Sciences}
\textbf{576} (1989), 500--535.
\newblock Author's publication list, entry 72:
\href{https://math.mit.edu/~rstan/pubs/}{math.mit.edu/\textasciitilde rstan/pubs}.
\end{thebibliography}
\end{document}
